% =====================================================================
%  chapters/minggu-10.tex
%  Minggu 10 — JavaScript Dasar: Variabel, Fungsi, Kondisi & DOM Selection
% =====================================================================

\chapter{JavaScript Dasar: Variabel, Fungsi, Kondisi \& DOM Selection}
\label{chap:mjs-10}

% ---------------------------------------------------------------------
% 1. CAPAIAN PEMBELAJARAN
% ---------------------------------------------------------------------
\begin{capaian}
Setelah menyelesaikan bab ini, mahasiswa diharapkan mampu:
\begin{itemize}
  \item menulis skrip JavaScript pertama dan menjalankannya di peramban;
  \item mendeklarasikan variabel dengan \texttt{let} dan \texttt{const}, serta membedakan tipe data dasar;
  \item mendefinisikan dan memanggil fungsi, termasuk \emph{arrow function};
  \item menggunakan percabangan \texttt{if}/\texttt{else} untuk pengambilan keputusan;
  \item memilih elemen HTML melalui \emph{DOM selection} (\texttt{getElementById}, \texttt{querySelector}, dsb.).
\end{itemize}
\end{capaian}

% ---------------------------------------------------------------------
% 2. TUJUAN PRAKTIK
% ---------------------------------------------------------------------
\begin{tujuanpraktik}
Pada sesi praktik ini, mahasiswa akan:
\begin{itemize}
  \item membuat file \texttt{.html} yang mengandung skrip JavaScript;
  \item berlatih menulis variabel, fungsi, dan kondisi;
  \item memanipulasi elemen HTML melalui \emph{DOM selection}.
\end{itemize}
\end{tujuanpraktik}

% =====================================================================
\section{Persiapan}
\label{sec:10-persiapan}
% =====================================================================

Yang Anda butuhkan:
\begin{itemize}
  \item VSCode (sudah terinstal dari minggu sebelumnya).
  \item File HTML dari minggu-minggu sebelumnya (atau buat file baru \texttt{index.html}).
  \item Peramban (Chrome atau Firefox).
\end{itemize}

\begin{catatan}
Untuk modul ini kita \textbf{TIDAK} membutuhkan Node.js. Kita menulis JavaScript langsung di dalam file HTML.
\end{catatan}

% =====================================================================
\section{Menulis Skrip Pertama}
\label{sec:10-srip-pertama}
% =====================================================================

\begin{langkah}
  \item Buat file baru bernama \texttt{minggu10.html} di folder proyek Anda.
  \item Isi dengan kode berikut.
  \item Buka file ini di peramban (klik kanan di VSCode $\rightarrow$ \textbf{Open with Live Server}).
  \item Buka \textbf{Console} peramban: tekan \texttt{F12} $\rightarrow$ tab \textbf{Console}.
\end{langkah}

\begin{lstlisting}[language=HTML, caption={Skrip JavaScript pertama}, label={lst:10-srip-pertama}]
<!DOCTYPE html>
<html lang="id">
<head>
    <meta charset="UTF-8">
    <meta name="viewport" content="width=device-width, initial-scale=1.0">
    <title>Minggu 10 -- JavaScript Pertama</title>
</head>
<body>
    <h1 id="judul">Halo, Dunia!</h1>
    <p id="paragraf">Ini adalah paragraf pertama saya.</p>

    <script>
        console.log("JavaScript berhasil dimuat!");
        document.getElementById("judul").style.color = "blue";
    </script>
</body>
</html>
\end{lstlisting}

\begin{tangkapanlayar}
  {assets/images/minggu-10/console-srip-pertama.png}
  {Console peramban menampilkan pesan ``JavaScript berhasil dimuat!''}
  {Buka file \texttt{minggu10.html} di peramban, lalu buka Console (\texttt{F12}) dan bandingkan hasilnya.}
\end{tangkapanlayar}

\subsection{Yang Terjadi}
\begin{itemize}
  \item \texttt{<script>} di dalam HTML adalah tempat menulis kode JavaScript.
  \item \texttt{console.log(...)} mencetak pesan ke Console peramban---sangat berguna untuk \emph{debugging}.
  \item \texttt{document.getElementById("judul")} mencari elemen HTML dengan \texttt{id="judul"}.
  \item \texttt{.style.color = "blue"} mengubah warna teks menjadi biru.
\end{itemize}

% =====================================================================
\section{Variabel \& Tipe Data}
\label{sec:10-variabel}
% =====================================================================

\subsection{Deklarasi Variabel}

Variabel adalah \textbf{wadah} untuk menyimpan data. Di JavaScript modern ada tiga cara mendeklarasikan variabel:

\begin{table}[htbp]
\centering
\caption{Perbandingan kata kunci deklarasi variabel}
\label{tab:10-variabel}
\begin{tabular}{lll}
\toprule
\textbf{Kata Kunci} & \textbf{Bisa Diubah?} & \textbf{Kapan Dipakai} \\
\midrule
\texttt{let}    & Ya   & Nilai yang mungkin berubah \\
\texttt{const}  & Tidak & Nilai yang tetap (konstanta) \\
\texttt{var}    & Ya   & HINDARI---cara lama, bermasalah pada \emph{scope} \\
\bottomrule
\end{tabular}
\end{table}

\begin{tip}
Selalu pakai \texttt{const}. Ganti ke \texttt{let} hanya jika nilai perlu diubah. Jangan pakai \texttt{var}.
\end{tip}

\begin{langkah}
  \item Tambahkan script berikut \textbf{sebelum} tag \texttt{</script>} yang sudah ada.
  \item \emph{Refresh} halaman dan lihat hasilnya di Console.
\end{langkah}

\begin{lstlisting}[language=JavaScript, caption={Deklarasi variabel dan perubahan nilai}, label={lst:10-variabel}]
// Variabel dengan const (tidak bisa diubah)
const nama = "Budi";
const umur = 19;

// Variabel dengan let (bisa diubah)
let IPK = 3.5;

console.log("Nama:", nama);
console.log("Umur:", umur);
console.log("IPK:", IPK);

// Mengubah nilai let
IPK = 3.8;
console.log("IPK setelah update:", IPK);
\end{lstlisting}

\begin{tangkapanlayar}
  {assets/images/minggu-10/console-variabel.png}
  {Console menampilkan output nama, umur, IPK, dan IPK setelah \emph{update}.}
  {Bandingkan output Console Anda dengan contoh di atas.}
\end{tangkapanlayar}

\subsection{Tipe Data Dasar}

\begin{lstlisting}[language=JavaScript, caption={Tipe data dasar JavaScript}, label={lst:10-tipedata}]
// String (teks) -- pakai tanda kutip
const pesan = "Halo Dunia";

// Number (angka) -- tanpa tanda kutip
const nilai = 100;

// Boolean (benar/salah)
const lulus = true;

// Array (daftar)
const buku = ["HTML", "CSS", "JavaScript"];

// Object (kumpulan properti)
const mahasiswa = {
    nama: "Budi",
    nim: "12345",
    jurusan: "Informatika"
};

console.log(pesan);           // "Halo Dunia"
console.log(nilai);           // 100
console.log(lulus);           // true
console.log(buku[0]);         // "HTML" (indeks mulai dari 0)
console.log(mahasiswa.nama);  // "Budi"
\end{lstlisting}

% =====================================================================
\section{Fungsi}
\label{sec:10-fungsi}
% =====================================================================

Fungsi adalah \textbf{blok kode yang bisa dipanggil berulang kali}. Bayangkan seperti resep masakan---Anda menulis resepnya sekali, lalu bisa dipakai kapan saja.

\begin{langkah}
  \item Tambahkan kode berikut ke dalam tag \texttt{<script>}.
  \item \emph{Refresh} halaman, lihat hasil di Console.
\end{langkah}

\begin{lstlisting}[language=JavaScript, caption={Fungsi deklarasi, nilai kembali, dan \emph{arrow function}}, label={lst:10-fungsi}]
// Fungsi sederhana -- deklarasi
function sapa(nama) {
    console.log("Halo, " + nama + "! Selamat belajar JavaScript.");
}

// Memanggil fungsi
sapa("Budi");
sapa("Ani");
sapa("Andi");

// Fungsi dengan nilai kembali (return)
function hitungLuas(panjang, lebar) {
    const luas = panjang * lebar;
    return luas;
}

const hasil = hitungLuas(10, 5);
console.log("Luas persegi panjang:", hasil);

// Fungsi panah (arrow function) -- cara modern
const kalikan = (a, b) => {
    return a * b;
};

console.log("3 x 4 =", kalikan(3, 4));
\end{lstlisting}

\begin{tangkapanlayar}
  {assets/images/minggu-10/console-fungsi.png}
  {Console menampilkan output pemanggilan fungsi \texttt{sapa()}, \texttt{hitungLuas()}, dan \texttt{kalikan()}.}
  {Bandingkan output Console Anda dengan contoh di atas.}
\end{tangkapanlayar}

\begin{itemize}
  \item \texttt{function sapa(nama)} mendefinisikan fungsi bernama \texttt{sapa} yang menerima satu parameter \texttt{nama}.
  \item \texttt{sapa("Budi")} memanggil fungsi dengan nilai \texttt{"Budi"} untuk parameter \texttt{nama}.
  \item \texttt{return luas} mengembalikan nilai hasil perhitungan agar bisa digunakan di luar fungsi.
  \item \emph{Arrow function} \texttt{=>} adalah cara menulis fungsi yang lebih singkat.
\end{itemize}

% =====================================================================
\section{Kondisi (If/Else)}
\label{sec:10-kondisi}
% =====================================================================

Kondisi membuat program bisa \textbf{mengambil keputusan}---seperti ``Jika hujan, bawa payung; jika tidak, bawa topi.''

\begin{langkah}
  \item Tambahkan kode berikut ke dalam tag \texttt{<script>}.
  \item \emph{Refresh} dan lihat output di Console.
\end{langkah}

\begin{lstlisting}[language=JavaScript, caption={Percabangan if/else dan operator perbandingan}, label={lst:10-kondisi}]
// Kondisi sederhana
const suhu = 32;

if (suhu > 30) {
    console.log("Panas! Minum air yang banyak.");
} else if (suhu >= 20) {
    console.log("Nyaman. Cocok untuk belajar.");
} else {
    console.log("Dingin! Pakai jaket.");
}

// Kondisi dengan string
const namaMahasiswa = "Budi";

if (namaMahasiswa === "Budi") {
    console.log("Halo Budi, selamat datang!");
} else {
    console.log("Siapa kamu?");
}

// Operator perbandingan
// ===  sama dengan (nilai DAN tipe sama)
// !==  tidak sama dengan
// >    lebih besar dari
// <    lebih kecil dari
// >=   lebih besar atau sama dengan
// <=   lebih kecil atau sama dengan

console.log(5 === 5);    // true
console.log(5 === "5");  // false (angka vs string)
console.log(5 == "5");   // true (hati-hati, bisa menyesatkan)
\end{lstlisting}

\begin{tangkapanlayar}
  {assets/images/minggu-10/console-kondisi.png}
  {Console menampilkan output kondisi if/else dan perbandingan.}
  {Bandingkan output Console Anda dengan contoh di atas.}
\end{tangkapanlayar}

% =====================================================================
\section{DOM Selection}
\label{sec:10-dom}
% =====================================================================

\textbf{DOM (Document Object Model)} adalah representasi struktur HTML sebagai ``pohon'' yang bisa diakses oleh JavaScript. \emph{DOM Selection} adalah cara JavaScript \textbf{mencari elemen} di halaman HTML.

\begin{langkah}
  \item Ganti isi file \texttt{minggu10.html} dengan versi lengkap berikut.
  \item Buka dengan Live Server dan lihat hasilnya.
  \item Buka Console (\texttt{F12} $\rightarrow$ Console) untuk melihat semua \emph{log}.
\end{langkah}

\begin{lstlisting}[language=HTML, caption={Demonstrasi DOM Selection}, label={lst:10-dom}]
<!DOCTYPE html>
<html lang="id">
<head>
    <meta charset="UTF-8">
    <meta name="viewport" content="width=device-width, initial-scale=1.0">
    <title>Minggu 10 -- DOM Selection</title>
    <style>
        body { font-family: Arial, sans-serif; margin: 20px; }
        .kotak { border: 2px solid #333; padding: 10px;
                 margin: 10px 0; border-radius: 5px; }
        .highlight { background-color: yellow; }
    </style>
</head>
<body>
    <h1 id="judul">Belajar DOM Selection</h1>
    <p class="teks">Paragraf pertama -- class="teks"</p>
    <p class="teks">Paragraf kedua -- class="teks"</p>
    <p id="paragraf-khusus">Paragraf khusus dengan ID</p>

    <div class="kotak" id="kotak1">Kotak 1</div>
    <div class="kotak" id="kotak2">Kotak 2</div>
    <div class="kotak" id="kotak3">Kotak 3</div>

    <ul id="daftar-mahasiswa">
        <li>Budi</li>
        <li>Ani</li>
        <li>Andi</li>
        <li>Sari</li>
    </ul>

    <script>
        // Cara 1: getElementById
        const judul = document.getElementById("judul");
        console.log("Judul:", judul.textContent);
        judul.style.color = "purple";

        // Cara 2: querySelector
        const paragrafKhusus =
            document.querySelector("#paragraf-khusus");
        console.log("Paragraf khusus:", paragrafKhusus.textContent);

        // Cara 3: querySelectorAll
        const semuaTeks = document.querySelectorAll(".teks");
        console.log("Jumlah paragraf .teks:", semuaTeks.length);
        semuaTeks.forEach((el, index) => {
            console.log(`  Teks ${index + 1}:`, el.textContent);
        });

        // Cara 4: getElementsByClassName
        const semuaKotak =
            document.getElementsByClassName("kotak");
        console.log("Jumlah kotak:", semuaKotak.length);
        for (let i = 0; i < semuaKotak.length; i++) {
            semuaKotak[i].style.backgroundColor = "lightblue";
            semuaKotak[i].style.padding = "15px";
        }

        // Cara 5: selector campuran
        const itemPertama = document.querySelector(
            "#daftar-mahasiswa li:first-child"
        );
        console.log("Mahasiswa pertama:", itemPertama.textContent);
        itemPertama.style.fontWeight = "bold";
        itemPertama.style.color = "red";
    </script>
</body>
</html>
\end{lstlisting}

\begin{tangkapanlayar}
  {assets/images/minggu-10/dom-selection-halaman.png}
  {Halaman web setelah dieksekusi: judul ungu, kotak biru, item pertama daftar mahasiswa berwarna merah tebal.}
  {Buka \texttt{minggu10.html} di peramban dan bandingkan tampilan Anda.}
\end{tangkapanlayar}

\begin{tangkapanlayar}
  {assets/images/minggu-10/dom-selection-console.png}
  {Console menampilkan semua \emph{log} DOM selection.}
  {Buka Console (\texttt{F12}) dan pastikan semua \emph{log} muncul.}
\end{tangkapanlayar}

\subsection{Ringkasan Metode DOM Selection}

\begin{table}[htbp]
\centering
\caption{Metode DOM Selection}
\label{tab:10-dom}
\begin{tabular}{lll}
\toprule
\textbf{Metode} & \textbf{Mencari} & \textbf{Mengembalikan} \\
\midrule
\texttt{getElementById("id")} & 1 elemen & Element \\
\texttt{querySelector("sel")} & 1 elemen (pertama cocok) & Element \\
\texttt{querySelectorAll("sel")} & Semua elemen & NodeList \\
\texttt{getElementsByClassName("cls")} & Semua elemen & HTMLCollection \\
\bottomrule
\end{tabular}
\end{table}

\begin{tip}
\texttt{querySelector} dan \texttt{querySelectorAll} paling fleksibel---bisa pakai semua \emph{selector} CSS (\texttt{id}, \texttt{class}, \texttt{tag}, \texttt{attribute}, bahkan \texttt{:nth-child}).
\end{tip}

% =====================================================================
\section{Latihan}
\label{sec:10-latihan}
% =====================================================================

\begin{latihan}[Variabel dan Fungsi]
Buat file baru \texttt{latihan10-1.html} dan buatlah:
\begin{enumerate}
  \item Variabel \texttt{const} untuk menyimpan nama lengkap Anda dan NIM.
  \item Variabel \texttt{let} untuk menyimpan nilai ujian (angka).
  \item Fungsi \texttt{beriNilai(skor)} yang mencetak ke Console:
  \begin{itemize}
    \item Jika skor $\geq$ 85: \texttt{"Excellent! Nilai Anda: [skor]"}
    \item Jika skor $\geq$ 70: \texttt{"Bagus! Nilai Anda: [skor]"}
    \item Jika skor $<$ 70: \texttt{"Perlu belajar lagi. Nilai Anda: [skor]"}
  \end{itemize}
  \item Panggil fungsi tersebut dengan tiga nilai berbeda.
\end{enumerate}

\textbf{Output yang diharapkan di Console:}
\begin{lstlisting}[language=JavaScript]
Excellent! Nilai Anda: 90
Bagus! Nilai Anda: 75
Perlu belajar lagi. Nilai Anda: 60
\end{lstlisting}

\begin{tangkapanlayar}
  {assets/images/minggu-10/latihan1-expected.png}
  {Console menampilkan output pencarian nilai sesuai skor.}
  {Jalankan file \texttt{latihan10-1.html} di peramban dan pastikan output sesuai.}
\end{tangkapanlayar}
\end{latihan}

\begin{latihan}[DOM Selection]
Buat file baru \texttt{latihan10-2.html} dengan HTML berikut:

\begin{lstlisting}[language=HTML, label={lst:10-latihan2-html}]
<!DOCTYPE html>
<html lang="id">
<head>
    <meta charset="UTF-8">
    <meta name="viewport" content="width=device-width, initial-scale=1.0">
    <title>Latihan 10-2</title>
</head>
<body>
    <h1 id="sapa">Selamat Datang</h1>
    <p class="info">Mata kuliah: Desain Web</p>
    <p class="info">Dosen: [Nama Dosen]</p>
    <p class="info">Hari: [Hari Kuliah]</p>
    <div class="kartu">Kartu Mahasiswa</div>
    <div class="kartu">Kartu UKT</div>
    <div class="kartu">Kartu Perpustakaan</div>
</body>
</html>
\end{lstlisting}

Tugas JavaScript:
\begin{enumerate}
  \item Ambil elemen \texttt{\#sapa} dan ubah teksnya menjadi nama Anda.
  \item Ambil semua elemen \texttt{.info} dan ubah warna setiap paragraf menjadi biru.
  \item Ambil semua elemen \texttt{.kartu} dan tambahkan \texttt{border 1px solid gray} pada setiap kartu.
\end{enumerate}

\begin{tangkapanlayar}
  {assets/images/minggu-10/latihan2-expected.png}
  {Halaman setelah JavaScript berjalan---judul berganti nama, paragraf info biru, kartu ada border.}
  {Jalankan file \texttt{latihan10-2.html} di peramban dan pastikan tampilan sesuai.}
\end{tangkapanlayar}
\end{latihan}

\begin{latihan}[Tantangan --- Opsional]
Buat file \texttt{latihan10-3.html} yang berisi:
\begin{itemize}
  \item Sebuah \emph{array} berisi minimal 5 nama mahasiswa.
  \item Fungsi \texttt{cariMahasiswa(namaArray, namaDicari)} yang mencetak ke Console:
  \begin{itemize}
    \item \texttt{"[nama] ditemukan di indeks [indeks]!"} jika ketemu.
    \item \texttt{"[nama] tidak ditemukan di daftar."} jika tidak ketemu.
  \end{itemize}
  \item Panggil fungsi dengan nama yang ada di \emph{array} dan nama yang tidak ada.
\end{itemize}

\textbf{Output yang diharapkan di Console:}
\begin{lstlisting}[language=JavaScript]
Budi ditemukan di indeks 0!
Siti tidak ditemukan di daftar.
\end{lstlisting}

\begin{tangkapanlayar}
  {assets/images/minggu-10/latihan3-expected.png}
  {Console menampilkan hasil pencarian nama dalam \emph{array}.}
  {Jalankan file \texttt{latihan10-3.html} di peramban.}
\end{tangkapanlayar}
\end{latihan}

% =====================================================================
\section{Troubleshooting}
\label{sec:10-troubleshooting}
% =====================================================================

\begin{peringatan}
\textbf{JavaScript tidak berjalan / Console kosong} --- pastikan \texttt{<script>} berada \textbf{sebelum} \texttt{</body>}, perhatikan huruf besar-kecil nama fungsi/variabel, dan tekan \texttt{Ctrl+S} sebelum \emph{refresh}.
\end{peringatan}

\begin{catatan}
\textbf{``Uncaught TypeError: Cannot read properties of null''} --- elemen belum ada di DOM saat script dijalankan. Pindahkan script ke akhir \texttt{<body>} atau bungkus dengan \texttt{DOMContentLoaded}.
\end{catatan}

\begin{lstlisting}[language=JavaScript, caption={Penanganan elemen belum tersedia}, label={lst:10-trouble}]
// Salah -- elemen belum ada di DOM
document.getElementById("judul").style.color = "red";

// Benar -- pindahkan script ke akhir body
// atau pakai DOMContentLoaded
document.addEventListener("DOMContentLoaded", function() {
    document.getElementById("judul").style.color = "red";
});
\end{lstlisting}

\begin{tip}
\textbf{Selector tidak menemukan elemen} --- buka DevTools $\rightarrow$ tab Elements $\rightarrow$ periksa apakah elemen yang dicari benar-benar ada di DOM. Pastikan \texttt{\#} untuk ID dan \texttt{.} untuk \emph{class}.
\end{tip}
